%==============================================================================
\section{Resultados}
\label{sec:results}

% --- Preenchido pelo pipeline de análise (05_Execution_Plan.md, Tier 1 primeiro).
%     Números vêm de generated/values.tex; macros [TODO] marcam execuções pendentes.

\subsection{Informação mútua vs.\ $R^2$}
\label{sec:res-mi}
A \Cref{tab:mi_vs_r2} coloca as medidas de informação ao lado da referência
linear $R^2$ na campanha PID, todas computadas sobre a série \emph{bruta}
pós-\emph{soak} (não detrendizada), para pareamento direto com o $R^2$ do
trabalho anterior (\cref{sec:res-correction} trata a MI \emph{reversível},
detrendizada, separadamente). A estimativa com \emph{binning}
(Miller--Madow) --- o estimador-manchete deste trabalho, por ser o único com
intervalo de confiança reportado --- dá $\MI{\slk}{\Tdie,\Vcore}=\valMIstv$~bits,
com a temperatura dominando ($\MI{\slk}{\Tdie}=\valMIst$ vs.\
$\MI{\slk}{\Vcore}=\valMIsv$~bits) e intervalo de 95\,\% por
\emph{moving-block bootstrap} $[\valMIciLo,\valMIciHi]$~bits. O estimador
KSG, aplicado como verificação independente sobre uma subamostra decimada,
produz $\valKSGstv$~bits, com seu próprio intervalo de 95\,\% por
\emph{bootstrap} em blocos $[\valKSGciLo,\valKSGciHi]$~bits --- os dois
intervalos se sobrepõem apenas estreitamente (em
$[\valKSGciLo,\valMIciHi]$), com os pontos estimados em extremos opostos da
sobreposição: uma diferença sistemática entre estimadores caracterizada com
barras de erro próprias abaixo, não uma concordância numérica.

\paragraph{O resultado de não-linearidade.} A comparação correta,
maçã-com-maçã, \emph{não} é a fração de entropia contra a fração de
variância (escalas diferentes), mas a MI medida contra a MI
\emph{Gaussiana-equivalente} implicada pelo ajuste linear,
$-\tfrac12\log_2(1-R^2)=\valMIlinEquiv$~bits para $R^2=\valRsqLinear$. A
dependência medida com \emph{binning} ($\valMIstv$~bits) \textbf{excede} esse
valor linear-equivalente em $\valMIexcessPctBinned\,\%$; o estimador KSG,
sobre a mesma comparação, dá um excesso de $\valMIexcessPctKSG\,\%$. Os dois
estimadores concordam na conclusão qualitativa --- existe informação
ambiental genuína que um modelo linear não consegue ver --- mas não no valor
exato do excesso; reportamos, portanto, uma faixa de
\textbf{$\valMIexcessPctBinned$--$\valMIexcessPctKSG\,\%$} conforme o
estimador, em vez de um único número pontual, como o suporte quantitativo
para substituir $R^2$ por informação mútua neste domínio. Essa faixa é o
intervalo entre as estimativas pontuais dos dois estimadores; propagando o
IC de cada um, o envelope completo é
$\valMIexcessCiLoBinned$--$\valMIexcessCiHiKSG\,\%$, cujo extremo inferior
permanece acima do equivalente linear.

\paragraph{Sensibilidade ao número de \emph{bins}.} $\Tdie$ e $\Vcore$ são
contínuos e foram discretizados em $\valMIbins$ \emph{bins} equiprováveis
para a estimativa acima (\cref{sec:th-diffent}). Repetindo a estimativa com
\emph{binning} para $12$, $24$, $32$ \emph{bins}, $\MI{\slk}{\Tdie,\Vcore}$ varia
em $[\valMIbinSensLo,\valMIbinSensHi]$~bits --- uma faixa de
$\sim$0.06~bits, da mesma ordem da largura do intervalo de confiança por
bootstrap. Mesmo no extremo inferior dessa faixa ($\valMIbinSensLo$~bits), o
valor permanece muito acima do piso nulo \emph{surrogate}
(\valMInullFloor~bits, abaixo) e acima do equivalente Gaussiano linear
($\valMIlinEquiv$~bits): nenhuma conclusão deste trabalho depende da escolha
particular de $\valMIbins$ \emph{bins}.

\paragraph{O excesso sobrevive à remoção da tendência? (achado negativo,
reportado honestamente).} O resultado de não-linearidade acima é
deliberadamente medido sobre a série \emph{bruta}, pareada com o $R^2$
também bruto do trabalho anterior. Isso deixa em aberto uma pergunta mais
forte: o excesso está no acoplamento reversível \emph{instantâneo}
$\Tdie,\Vcore\to\slk$, ou concentrado na tendência lenta que envelhecimento
e $\Tdie,\Vcore$ compartilham ao longo de uma campanha de
\valDurationHours~h? Repetindo a mesma comparação sobre as séries
detrendizadas usadas na correção PVT (\cref{sec:res-correction}):
$R^2_{\text{detr}}=\valRsqDetrended$, equivalente Gaussiano
$=\valMIlinEquivDetrended$~bits; a MI medida com \emph{binning} é
$\valMIdtBinned$~bits (IC 95\,\% $[\valMIdtCiLo,\valMIdtCiHi]$) e o KSG,
aplicado à mesma série detrendizada com seu próprio \emph{bootstrap} em
blocos, dá $\valMIdtKSG$~bits (IC 95\,\% $[\valMIdtKsgCiLo,\valMIdtKsgCiHi]$)
--- \textbf{os dois estimadores ficam no equivalente Gaussiano linear ou
abaixo dele} (déficit de $\valMIexcessPctDetrended$--$\valMIexcessPctDetrKSG\,\%$
conforme o estimador; o extremo superior do IC do KSG, $\valMIdtKsgCiHi$,
toca o equivalente Gaussiano $\valMIlinEquivDetrended$, então o KSG é
compatível com paridade exata, enquanto o binned fica claramente abaixo).
Os dois valores detrendizados, binned ($\valMIdtBinned$) e KSG
($\valMIdtKSG$), diferem entre si pela mesma razão que os valores brutos da
\cref{sec:res-mi} diferem (\valMIstv\ vs.\ \valKSGstv): é a mesma série
detrendizada em ambos os casos, apenas estimada por famílias de estimador
distintas (\emph{binning} com correção de Miller--Madow vs.\
$k$-vizinhos-mais-próximos) --- não um erro de pré-processamento nem duas
séries diferentes.

\emph{Como um déficit numérico deve ser lido.} Isto é reportado
honestamente, no mesmo espírito dos demais resultados negativos deste
trabalho (\cref{sec:res-correction,sec:res-wiener}): o excesso não-linear de
$\valMIexcessPctBinned$--$\valMIexcessPctKSG\,\%$ medido na série bruta
\emph{não sobrevive} à remoção da tendência lenta compartilhada, e portanto
não pode ser atribuído com confiança a uma não-linearidade
\emph{instantânea} do tipo Arrhenius/super-linear na relação reversível
$\Tdie,\Vcore\to\slk$ --- ele está concentrado, ao menos em parte, na
interação entre a deriva de envelhecimento e a deriva lenta de
$\Tdie,\Vcore$ ao longo da campanha. Esse déficit numérico, porém,
\textbf{não} deve ser lido como ``dependência menor que a linear'': o
equivalente Gaussiano é um construto de variáveis contínuas, e para um
resíduo quantizado em poucas contagens a MI discreta pode legitimamente
ficar abaixo dele (o teto entrópico do resíduo e o viés
descendente do estimador nesse regime de MI pequena e $\neff$ reduzido
empurram nessa direção); a leitura defensável é que \textbf{nenhum excesso
não-linear é resolvível acima do piso do estimador} nesta série, não que a
relação reversível seja ``sub-linear''. A comparação bruta permanece o
pareamento correto e honesto com o $R^2$ publicado do trabalho anterior
(também medido sobre série bruta, sem \emph{detrend}), mas a leitura causal
--- ``a física é não-linear, por isso a MI excede o $R^2$'' --- deve ser
lida como a hipótese que motivou este trabalho, não como algo que este
resultado demonstra; a \cref{sec:discussion} retoma este ponto.

\paragraph{Piso nulo \emph{surrogate}.} Como estimadores de MI são viesados
para cima em amostras finitas e autocorrelacionadas, o valor medido precisa
ser verificado contra o piso que o estimador produz sob independência
genuína. Deslocando circularmente o fluxo $(\Tdie,\Vcore)$ contra $\slk$ em
defasagens aleatórias grandes --- destruindo a dependência cruzada enquanto
preserva a marginal \emph{e} a autocorrelação de cada canal --- obtém-se um
$\MI{\slk}{\Tdie,\Vcore}$ nulo de média $\valMInullMean$~bits e percentil 95
de $\valMInullFloor$~bits. O valor medido de $\valMIstv$~bits é
$\approx\valMIfloorRatio\times$ esse piso, de modo que a dependência é
esmagadoramente real, e não viés de estimador. O mesmo raciocínio se aplica
ao excesso sobre o equivalente Gaussiano linear da série bruta: descontando
o viés médio do estimador ($\valMInullMean$~bits) do excesso com
\emph{binning} ($\valMIexcessPctBinned\,\%$), sobra $\valMIexcessFloorAdjBits$~bits
($\approx\valMIexcessFloorAdjPct\,\%$) ainda positivo --- o excesso da série
bruta não é, ele próprio, um artefato do piso de viés (isso não contradiz o
achado do parágrafo anterior: o piso de viés e a confusão com a tendência
compartilhada são duas ameaças distintas, e o excesso bruto sobrevive à
primeira mas não à segunda). O mesmo piso enquadra o
resíduo reversível pós-correção ($\valResidMI$~bits,
\cref{sec:res-correction}): ter reduzido o acoplamento remanescente à mesma
ordem de grandeza do piso nulo é exatamente o que ``o ambiente foi removido''
significa em termos de informação.

\begin{figure}[t]
  \centering
  \includegraphics[width=0.86\linewidth]{fig_mi.pdf}
  \caption{Informação mútua por canal e conjunta entre \emph{slack} e o
  ambiente (bits), estimador com \emph{binning}. A conjunta
  $\MI{\slk}{\Tdie,\Vcore}$ excede a Gaussiana-equivalente do $R^2$ linear
  (tracejado) em $\valMIexcessPctBinned\%$ (KSG: $\valMIexcessPctKSG\%$), e
  fica bem acima do piso nulo \emph{surrogate} (pontilhado), confirmando
  acoplamento genuíno (a atribuição do excesso à não-linearidade é
  discutida na \cref{sec:discussion}).}
  \label{fig:mi}
\end{figure}

% Tabela principal: medidas de informação vs. a referência linear R^2.
\begin{table}[t]
  \centering
  \caption{Medidas de informação (bits) vs.\ a referência linear $R^2$ na
           campanha PID. A $\MI{\slk}{\Tdie,\Vcore}$ medida excede a MI
           Gaussiana-equivalente do ajuste linear
           ($-\tfrac12\log_2(1-R^2)$) em \valMIexcessPct\,\% --- um excesso
           real (sobrevive ao piso de viés do estimador), cuja atribuição à
           não-linearidade instantânea, porém, não sobrevive à remoção da
           tendência compartilhada (\cref{sec:res-mi,sec:discussion}). As
           colunas de trabalho anterior são as estatísticas publicadas de
           fidelidade térmica.}
  \label{tab:mi_vs_r2}
  \begin{tabular}{lccc}
    \toprule
     & PID (nova execução) & PID (anterior) & Bang-bang (anterior) \\
    \midrule
    $R^2$ linear $(\slk\mid\Tdie,\Vcore)$       & \valRsqLinear & \valRsqPid & \valRsqBangBang \\
    \;$\hookrightarrow$ MI Gaussiana-equiv.\ (bits) & \valMIlinEquiv & --- & --- \\
    $\Ent{\slk}$ (bits)                      & \valHs        & ---        & --- \\
    $\Ent{\slk\mid\Tdie,\Vcore}$ (bits)      & \valHscondTV  & ---        & --- \\
    $\MI{\slk}{\Tdie}$ (bits)                & \valMIst      & ---        & --- \\
    $\MI{\slk}{\Vcore}$ (bits)               & \valMIsv      & ---        & --- \\
    $\MI{\slk}{\Tdie,\Vcore}$ com \emph{binning} (bits)  & \valMIstv     & ---        & --- \\
    \;$\hookrightarrow$ verificação cruzada KSG (bits) & \valKSGstv  & ---        & --- \\
    \;$\hookrightarrow$ IC 95\,\% (bootstrap em blocos) & [\valMIciLo,\,\valMIciHi] & --- & --- \\
    fração de informação $\MI{}{}/\Ent{}$         & \valInfoFraction & ---     & --- \\
    \bottomrule
  \end{tabular}
\end{table}


\subsection{Correção PVT como recuperação de informação}
\label{sec:res-correction}
Ajustamos $\hat{\dpvt}(\Tdie,\Vcore)$ e formamos $\scorr=\slk-\hat{\dpvt}$. A
correção por mínimos quadrados ordinários (linear) zera o acoplamento linear
por construção ($\operatorname{corr}(\scorr,\Tdie)=+0.00$,
$\operatorname{corr}(\scorr,\Vcore)=-0.00$); a pergunta metodologicamente
relevante é quanta informação mútua \emph{reversível} ela remove. Medido
sobre sinais detrendizados --- para que a deriva lenta de envelhecimento não
se disfarce de acoplamento ambiental --- o acoplamento reversível cai de
$\valPreMIrev$~bits antes da correção para $\valResidMI$~bits depois, ou
seja, a correção linear remove \textbf{\valMIreductionPct\,\%} da informação
reversível ambiente--sensor, deixando um resíduo no piso de ruído do
estimador. O desvio-padrão residual reversível é $\valResidSigmaLin$~contagens,
consistente com o resíduo PID publicado no trabalho anterior
($\valResidPid$~contagens) --- uma verificação cruzada independente do
\emph{pipeline}.

A correção não-linear motivada pela física (tipo Arrhenius em $\Tdie$,
super-linear em $\Vcore$) remove marginalmente mais MI residual
($\valResidMInl$ vs.\ $\valResidMIlin$~bits), mas reduz a variância residual
de forma negligenciável; o critério MDL/BIC (\cref{eq:mdl},
\cref{sec:th-mdl}) portanto \textbf{seleciona o modelo linear}
($\valCorrModel$), julgando os parâmetros extras injustificados neste ponto
de estresse --- um veredito de comprimento de descrição revisitado na
\cref{sec:discussion}.

\paragraph{Recuperação do envelhecimento (resultado-chave).} Remover a
contribuição PVT não apenas limpa o sinal --- ela \emph{expõe} a trajetória
latente de envelhecimento. A correlação do \emph{slack} com o tempo se
fortalece de $\operatorname{corr}(\slk,\tm)=\valCorrStRaw$ no fluxo bruto
para $\operatorname{corr}(\scorr,\tm)=\valCorrStCorr$ após a correção:
retirar a informação ambiental reversível torna a deriva permanente de
envelhecimento a estrutura dominante no resíduo, precisamente o objetivo de
recuperação de informação da \cref{eq:model} e a âncora empírica para o
limite de RUL da \cref{sec:res-rul}.

\begin{figure}[t]
  \centering
  \includegraphics[width=\linewidth]{fig_correction.pdf}
  \caption{(a) O \emph{slack} bruto acompanha a temperatura do \emph{die}
  (corr $\valCorrST$). (b) Após a correção PVT o acoplamento reversível
  desaparece e a tendência permanente de envelhecimento é \emph{exposta}
  (a correlação \emph{slack}--tempo se fortalece de $\valCorrStRaw$ para
  $\valCorrStCorr$); a linha tracejada é a taxa de envelhecimento ajustada.}
  \label{fig:correction}
\end{figure}

\subsection{Distância distribucional entre bancadas}
\label{sec:res-kl}
Duas execuções estabilizadas de $\sim$24~h em bancadas diferentes --- a
configuração SBCCI não-PID e a referência controlada por PID --- permitem
colocar um número em ``quão longe o ambiente não-conforme empurra a
distribuição do \emph{slack}''. Ambas as execuções são no mesmo dispositivo
de teste dedicado, distinto do dispositivo de envelhecimento de 214~h das
seções anteriores; elas compartilham esse dispositivo de teste, de modo que
a divergência isola a diferença de bancada/ambiente e não a variação
dispositivo-a-dispositivo. As bancadas operam em pontos de operação
diferentes, então a comparação é feita sobre a flutuação centrada do
\emph{slack}, isolando o caráter do ruído do ponto de ajuste. O ruído de
\emph{slack} é de $\valSlackStdSBCCI$ contagens (SBCCI) vs.\
$\valSlackStdPID$ contagens (PID), acompanhando o ruído de temperatura
($\valTstdSBCCI$ vs.\ $\valTstdPID$~$^\circ$C) --- a contaminação é de
origem térmica.

A divergência de Kullback--Leibler é fortemente \emph{assimétrica}:
$\KL{p_{\text{SBCCI}}}{p_{\text{PID}}}=\valKLsp$~bits (o custo de descrever a
bancada ruidosa com a referência PID apertada) contra
$\KL{p_{\text{PID}}}{p_{\text{SBCCI}}}=\valKLps$~bits --- uma ilustração
direta do ponto da disciplina de que KL não é uma métrica. A divergência de
Jensen--Shannon, simétrica, é $\valJSbench$~bits (de um máximo de 1), uma
única ``distância entre ambientes de validação'' transferível. Uma forma
fechada Gaussiana de média nula dá $\valKLgauss$~bits para
$\KL{p_{\text{SBCCI}}}{p_{\text{PID}}}$; o valor empírico é maior
($\valKLsp$~bits), e esse excesso é, por si só, evidência de que a bancada
não-conforme injeta excursões \emph{não-Gaussianas}, de cauda pesada --- o
que motiva a análise de separação cega da \cref{sec:res-ica}.

\begin{figure}[t]
  \centering
  \includegraphics[width=0.86\linewidth]{fig_kl.pdf}
  \caption{Distribuições da flutuação centrada de \emph{slack} das duas
  bancadas (mesmo dispositivo de teste). A bancada SBCCI não-PID é
  marcadamente mais larga que a referência PID; a divergência
  $\KL{p_{\text{SBCCI}}}{p_{\text{PID}}}=\valKLsp$~bits (JS $=\valJSbench$~bits)
  quantifica a lacuna de fidelidade de bancada.}
  \label{fig:kl}
\end{figure}

\subsection{Resolução efetiva e confiabilidade de alarme}
\label{sec:res-capacity}
Tomando o ruído reversível pós-correção $\sigma_{\text{noise}}=\valSigNoise$
contagens e a dispersão do sinal de envelhecimento
$\sigma_{\text{sig}}=\valSigSignal$ contagens, a heurística de capacidade
(\cref{eq:capacity}, derivada na \cref{sec:th-capacity} a partir do
argumento de entropia máxima Gaussiana) dá $\mathrm{SNR}=\valSNR$,
$C=\valCapacityBits$~bits e, portanto, \valCapacityStates{} estados de
envelhecimento distinguíveis ao longo desta campanha --- um número
deliberadamente modesto, que reflete uma excursão de envelhecimento de
apenas $\sim\!\sqrt{\valSNR}\times$ o piso de ruído, e que a métrica
residual-$\sigma$ anterior não conseguia expressar.

A decisão de alarme de envelhecimento é modelada como um BSC cujo
cruzamento $p$ decorre do SNR de detecção por janela (\cref{eq:bsc}). A
confiabilidade é fortemente \emph{dependente da escala de tempo}: uma
decisão por hora é essencialmente um lançamento de moeda ($p\approx0.49$,
$\Cbsc\approx0$) porque uma hora de deriva está enterrada no ruído, enquanto
uma janela operacional de \valBSCwindow~h dá $p=\valBSCp$ e
$\Cbsc=\valBSCcap$~bits/decisão (quase perfeito), e a campanha completa é
efetivamente certa. A transição entre esses regimes coincide com o tempo
mínimo de observação derivado a seguir.

\begin{figure}[t]
  \centering
  \includegraphics[width=0.82\linewidth]{fig_capacity_rul.pdf}
  \caption{Capacidade do canal de alarme de envelhecimento
  $\Cbsc=1-\Hb(p)$ em função do comprimento da janela de decisão. A
  confiabilidade sobe de $\approx0$ (por hora, enterrada no ruído) para
  $\approx1$ bit, com a transição coincidindo com o tempo mínimo de
  observação de Cram\'er--Rao ($\approx\valMinObsTime$~h, tracejado) --- as
  três figuras de mérito são mutuamente consistentes.}
  \label{fig:capacity}
\end{figure}

\subsection{Limite de detectabilidade de RUL}
\label{sec:res-rul}
Colocando o RUL como estimação por mínimos quadrados da inclinação de
envelhecimento, e usando o ruído medido com o tamanho efetivo de amostra
$\neff=\valNeff$ ($\tauint=\valTauInt$~amostras, ou seja,
$\sim$5~min de decorrelação), o limite de Cram\'er--Rao
(\cref{sec:th-estimation}) dá um desvio-padrão da estimativa de inclinação
de $\valCRLBsd$~m-cnt/h ao longo da campanha de $\sim$214~h, portanto uma
taxa mínima detectável de envelhecimento (a $3\sigma$) de
\textbf{\valMinDetectRate~m-cnt/h}. A taxa de envelhecimento corrigida
medida é $\valAgingSlope$~m-cnt/h --- aproximadamente $15\times$ acima desse
piso, de modo que a tendência é firmemente detectável. Equivalentemente, o
tempo mínimo de observação para resolver a tendência medida a $3\sigma$ é
\textbf{$\sim$\valMinObsTime~h}, bem dentro desta campanha (e explicando por
que as execuções anteriores de $\sim$23~h não conseguiam separar o
envelhecimento do piso reversível). Esta é a afirmação honesta de
detectabilidade que o projeto prometeu: não uma predição de vida útil, mas
um limite limitado por ruído sobre \emph{quando} o envelhecimento se torna
observável.

\paragraph{Barras de credibilidade Bayesianas.} O limite frequentista é
complementado por uma regressão linear Bayesiana de $\scorr$ sobre o tempo,
ajustada sobre os dados decimados até os $\neff$ pontos efetivamente
independentes (para que a verossimilhança i.i.d.\ seja válida) sob uma
\emph{prior} de Jeffreys, dando uma posterior $t$ de Student para a
inclinação. A taxa de envelhecimento posterior é $\valBayesSlope$~m-cnt/h
com um intervalo de credibilidade de $95\%$ de
$[\valBayesSlopeLo,\valBayesSlopeHi]$~m-cnt/h, e a probabilidade posterior de
que a inclinação seja genuinamente negativa é de $\valBayesPaging\,\%$ ---
numericamente indistinguível de certeza; a leitura dessa deriva residual
\textbf{como envelhecimento} permanece condicionada à hipótese de
identificação deste estudo de dispositivo único (\cref{sec:discussion},
ameaças à validade). Propagar a posterior da inclinação pela relação de detectabilidade
dá um intervalo de credibilidade de $[\valBayesMotLo,\valBayesMotHi]$~h para
o tempo mínimo de observação, contendo a estimativa pontual de
$\valMinObsTime$~h. Os dois paradigmas concordam, e o tratamento Bayesiano
fornece as barras de erro calibradas que um futuro estudo de RUL propagaria
adiante.

\subsection{Uma trajetória de envelhecimento denoised e o RUL como distribuição (Wiener/Kalman)}
\label{sec:res-wiener}
Os resultados de Cram\'er--Rao e Bayesiano acima fornecem uma única taxa
para toda a campanha e sua barra de erro. Aqui fazemos, em vez disso, a
pergunta de estimação sequencial da \cref{sec:th-wiener}: fazendo a média em
blocos de $\scorr$ em $\valNBlocks$ observações quase-independentes (uma a
cada $2\tauint$, ou seja, uma unidade de $\neff$ cada), o ajuste por método
dos momentos da \cref{eq:wiener-model} dá uma difusão de Wiener
$\sigma_B=\valWienerSigmaB$ contagens/$\sqrt{\text{h}}$ e um ruído de
observação por média de bloco $R$ notavelmente \emph{abaixo} do
$\sigma_{\text{noise}}^2$ nativo por amostra (\cref{sec:res-capacity}) ---
evidência direta de que a média sobre uma janela completa de decorrelação
ainda remove estrutura sub-$\tauint$ que o tempo de autocorrelação único,
sozinho, não captura. O MLE do ruído de taxa colapsa para seu piso numérico
($\sigma_{\mathrm{rate}}^2\!\approx\!\valWienerQRate$): os dados preferem um
processo de Wiener de deriva \emph{constante} a um variável no tempo, uma
segunda confirmação independente --- ao lado do veredito de MDL da
\cref{sec:res-correction} --- de que esta campanha única, de ponto de
operação estreito, não contém informação de curvatura suficiente para
sustentar um modelo mais flexível. A taxa suavizada por Kalman/RTS ao final
da campanha é $\valWienerRateEnd$~m-cnt/h, consistente em sinal e ordem de
grandeza com a estimativa CRLB/Bayesiana da \cref{sec:res-rul}
($\sim\!-18$~m-cnt/h); a diferença de $\sim$15\% entre as duas taxas
derivadas independentemente reflete os estimadores distintos (OLS sobre toda
a série vs.\ um estado filtrado, com média em blocos), não uma contradição.
O desvio-padrão posterior do nível suavizado é essencialmente plano ao longo
da campanha ($\valWienerLevelSDstart\to\valWienerLevelSDend$ contagens): o
filtro atinge sua incerteza de regime permanente dentro dos primeiros
blocos e não continua encolhendo indefinidamente, exatamente como um filtro
linear-Gaussiano com $R$ e $\sigma_B^2$ constantes prevê.

\begin{figure}[t]
  \centering
  \includegraphics[width=\linewidth]{fig_wiener_rul.pdf}
  \caption{(a) Média de bloco de $\scorr$ (claro) e a trajetória de
  envelhecimento suavizada por Kalman/RTS $\hat\Xlat(t)$ com uma banda de
  credibilidade de 95\%. (b) A distribuição de RUL Gaussiana-inversa
  (\cref{eq:invgauss}) em dois limiares ancorados no Tier 1: um incremento
  do piso de ruído ($D=\sigma_{\text{noise}}$) e uma excursão em escala de
  sinal já observada ($D=\sigma_{\text{signal}}$).}
  \label{fig:wiener}
\end{figure}

\paragraph{RUL como distribuição.} Em $D=\sigma_{\text{noise}}=\valDnoise$
contagens (o mesmo piso de ruído que limita o sensor a
\valCapacityStates{} estados distinguíveis, \cref{sec:res-capacity}), o
tempo de primeira passagem de Wiener tem média $\valRULmeanNoise$~h mas
mediana $\valRULmedianNoise$~h; em $D=\sigma_{\text{signal}}=\valDsignal$
contagens, média $\valRULmeanSignal$~h e mediana $\valRULmedianSignal$~h. A
grande diferença entre média e mediana em ambos os casos não é um artefato:
a Gaussiana inversa é fortemente assimétrica à direita exatamente quando o
limiar $D$ é pequeno em relação à escala de difusão $\sigma_B$, que é
precisamente o regime de SNR baixo ($\mathrm{SNR}=\valSNR$,
\cref{sec:res-capacity}) medido aqui --- \emph{a mesma limitação de ruído
que restringe os estados resolvíveis do sensor também torna as predições de
RUL de curto horizonte altamente incertas}, uma ligação quantitativa entre o
resultado de capacidade e o resultado de RUL que uma única estimativa
pontual não consegue expressar. Isso é oferecido como uma
\emph{metodologia} para uma distribuição completa de RUL, exercitada em
limiares auto-referenciados porque, como em todo este trabalho, não existe
um limiar de falha calibrado externamente para este sensor; fornecer um
(\cref{sec:conclusion}) transforma a \cref{eq:invgauss} diretamente em uma
distribuição de RUL operacional.

\subsection{Separação cega de fontes (ICA)}
\label{sec:res-ica}
Como contraparte livre de modelo à correção por regressão da
\cref{sec:res-correction}, o FastICA \cite{hyvarinen2001} (contraste de
negentropia, \cref{eq:negentropy} na \cref{sec:th-ica}) foi aplicado aos três canais
observados $(\slk,\Tdie,\Vcore)$ isoladamente, sem nenhuma forma funcional
assumida para $\dpvt$. O componente recuperado que acompanha o tempo bate
com o resíduo de envelhecimento baseado em regressão $\scorr$ em
$|\mathrm{corr}|=\valICAagingScorr$ --- \emph{o ICA livre de modelo chega
independentemente a essencialmente o mesmo sinal de envelhecimento que a
regressão produziu}, o que valida cruzadamente a correção PVT sem usar o
acoplamento medido $\Tdie,\Vcore$. Um segundo componente se alinha com a
temperatura ($|\mathrm{corr}(\cdot,\Tdie)|=\valICAthermalT$). Ambas as
fontes recuperadas são fortemente não-Gaussianas (curtose em excesso
$\valICAkurtAging$ e $\valICAkurtThermal$), confirmando a premissa de que a
não-Gaussianidade é a característica que distingue as fontes físicas --- e
consistente com as excursões de cauda pesada que a análise KL
(\cref{sec:res-kl}) expôs.

A separação é \emph{parcial}, e é reportada como tal: a informação mútua
residual entre pares de componentes recuperados atinge $\valICAresidMI$~bits
(a independência ideal é $0$), porque a deriva lenta e não-estacionária de
envelhecimento e o ruído de cauda pesada satisfazem apenas
aproximadamente as hipóteses de i.i.d./estacionariedade do ICA. O ICA é,
portanto, corroborativo aqui --- uma rota independente, com poucas hipóteses,
para o mesmo componente de envelhecimento --- em vez de uma desmistura
limpa. O ICA recebe os três mesmos canais observados $(\slk,\Tdie,\Vcore)$ da
regressão e busca no mesmo espaço linear de 3 dimensões; a diferença está no
\emph{critério}, não no espaço de busca --- mínimos quadrados de um lado,
independência/não-Gaussianidade (negentropia) do outro. Que dois critérios
tão distintos, aplicados às mesmas observações, cheguem a componentes de
envelhecimento com $|\mathrm{corr}|=\valICAagingScorr$ (isto é,
$\approx\valICAagingVarExpl\,\%$ de variância compartilhada --- não
$\valICAagingVarExpl\,\%$ do ``sinal recuperado'', uma leitura que confundiria
correlação com fração de sinal) é, ainda assim, uma corroboração de critério
não trivial.

\begin{figure}[t]
  \centering
  \includegraphics[width=\linewidth]{fig_ica.pdf}
  \caption{O componente de envelhecimento recuperado pelo ICA (roxo)
  sobreposto ao resíduo baseado em regressão $\scorr$ (verde), ambos
  padronizados ($z$-score). O ICA livre de modelo, dado apenas
  $(\slk,\Tdie,\Vcore)$, acompanha o sinal de envelhecimento da regressão em
  $|\mathrm{corr}|=\valICAagingScorr$.}
  \label{fig:ica}
\end{figure}
